\bta{2006甲A}

\begin{enumerate}
	%\renewcommand{\labelenumi}{\arabic{enumi}.}
	% A(\Alph) a(\alph) I(\Roman) i(\roman) 1(\arabic)
	%设定全局标号series=example	%引用全局变量resume=example
	%[topsep=-0.3em,parsep=-0.3em,itemsep=-0.3em,partopsep=-0.3em]
	%可使用leftmargin调整列表环境左边的空白长度 [leftmargin=0em]
	\item
（$30'$）两个线性算符$\hat{A}$和$\hat{B}$满足下列关系：$\hat{A}^2=0,\quad \hat{A}\hat{A}^\dagger+\hat{A}^\dagger \hat{A}=1,\quad \hat{B}=\hat{A}^\dagger\hat{A}$。
\begin{enumerate}
	%\renewcommand{\labelenumi}{\arabic{enumi}.}
	% A(\Alph) a(\alph) I(\Roman) i(\roman) 1(\arabic)
	%设定全局标号series=example	%引用全局变量resume=example
	%[topsep=-0.3em,parsep=-0.3em,itemsep=-0.3em,partopsep=-0.3em]
	%可使用leftmargin调整列表环境左边的空白长度 [leftmargin=0em]
	\item
求证$\hat{B}^2=\hat{B}$；

\item 
求在$\hat{B}$表象中$\hat{A}$和$\hat{B}$的表达式。
	
\end{enumerate}


\banswer{
	\begin{enumerate}
		\item 由 $A^2=0$ 可知 $(A^\dagger)^2=0$。于是
		\begin{align*}
			B^2&=A^\dagger A A^\dagger A
			=A^\dagger(1-A^\dagger A)A\\
			&=A^\dagger A-(A^\dagger)^2A^2=A^\dagger A=B.
		\end{align*}

		\item 由 $B^2=B$，本征值满足 $b^2=b$，只能为 $0,1$。
		对归一化的 $|0\rangle$ 有 $\|A|0\rangle\|^2=\langle0|B|0\rangle=0$，
		所以 $A|0\rangle=0$。令 $|1\rangle=A^\dagger|0\rangle$，则由
		$AA^\dagger=I-B$ 得 $\langle1|1\rangle=1$、$A|1\rangle=|0\rangle$，
		并可验证 $B|1\rangle=|1\rangle$。
		在这一不可约的二维子空间中，按 $|0\rangle,|1\rangle$ 排列基矢，得到
		\[
			\boxed{B=\begin{pmatrix}0&0\\0&1\end{pmatrix},
			\qquad
			A=\begin{pmatrix}0&1\\0&0\end{pmatrix},
			\qquad
			A^\dagger=\begin{pmatrix}0&0\\1&0\end{pmatrix}}.
		\]
		若 $B=0$ 与 $B=1$ 的本征空间还具有额外简并，$A^\dagger$ 将两个空间的正交基一一配对，
		一般表示是这些 $2\times2$ 块的直和（或将每个矩阵元乘以简并空间的单位阵）。
		交换基矢顺序时，只需同时交换矩阵的行列。
	\end{enumerate}
}


\item 
（$30'$）粒子在势场$V(x)=A|x|^n\quad(-\infty <x<\infty,A>0)$中运动，试用不确定度关系估算基态能量。

\banswer{
	以 $d=\Delta x$ 表示基态的典型尺度，并取 $\langle x\rangle=\langle p\rangle=0$。由 $\Delta p\,d\geq\hbar/2$，估算能量为
	\[
		E(d)\simeq\frac{\hbar^2}{8m d^2}+A d^n.
	\]
	令 $dE/dd=0$，得到
	\[
		d_0=\left(\frac{\hbar^2}{4mnA}\right)^{1/(n+2)},
		\qquad
		\frac{\hbar^2}{8md_0^2}=\frac n2Ad_0^n.
	\]
	所以基态能量的量级估计为
	\[
		\boxed{E_0\simeq\frac{n+2}{2}A
		\left(\frac{\hbar^2}{4mnA}\right)^{n/(n+2)}}.
	\]
	测不准估算法只能确定这一尺度及相应的数值近似；精确系数一般需解 Schr\"odinger 方程。
}


\item 
（$30'$）设体系的哈密顿量$\hat{H}$从依赖于某一参量$\lambda$，又设体系处于某一束缚定态，其能量和本征函数分别记为$E_n$和$\psi_n(r)$。
\begin{enumerate}
	%\renewcommand{\labelenumi}{\arabic{enumi}.}
	% A(\Alph) a(\alph) I(\Roman) i(\roman) 1(\arabic)
	%设定全局标号series=example	%引用全局变量resume=example
	%[topsep=-0.3em,parsep=-0.3em,itemsep=-0.3em,partopsep=-0.3em]
	%可使用leftmargin调整列表环境左边的空白长度 [leftmargin=0em]
	\item
证明费曼——海尔曼定理：
$$\frac{\partial E_n}{\partial \lambda}=\int \psi^*_n(\vec{r})\frac{\partial \hat{H}}{\partial \lambda}\psi_n(\vec{r})d\vec{r}$$

\item 
利用费曼——海尔曼定理，求氢原子各束缚态的平均动能。\\
（提示：氢原子能级公式$E_n=-\frac{\mu e^4}{2\hbar^2}\cdot\frac{1}{n^2}$）
	
\end{enumerate}


\banswer{
	\begin{enumerate}
		\item 对归一化本征方程 $H(\lambda)|n\rangle=E_n|n\rangle$ 关于 $\lambda$ 求导，
		\[
			\frac{\partial H}{\partial\lambda}|n\rangle
			+H\frac{\partial|n\rangle}{\partial\lambda}
			=\frac{\partial E_n}{\partial\lambda}|n\rangle
			+E_n\frac{\partial|n\rangle}{\partial\lambda}.
		\]
		左乘 $\langle n|$，利用 $\langle n|H=E_n\langle n|$，含 $\partial|n\rangle/\partial\lambda$ 的两项抵消，故
		\[
			\boxed{\frac{\partial E_n}{\partial\lambda}
			=\left\langle n\left|\frac{\partial H}{\partial\lambda}\right|n\right\rangle}
			=\int\psi_n^*\frac{\partial H}{\partial\lambda}\psi_n\,d^3r.
		\]

		\item 取参数 $\lambda=\mu$。氢原子哈密顿量与能量为
		\[
			H=\frac{p^2}{2\mu}-\frac{e^2}{r},\qquad
			E_n=-\frac{\mu e^4}{2\hbar^2n^2}.
		\]
		于是
		\[
			\frac{\partial H}{\partial\mu}=-\frac{p^2}{2\mu^2}=-\frac{T}{\mu},
			\qquad
			\frac{\partial E_n}{\partial\mu}=\frac{E_n}{\mu}.
		\]
		由费曼--海尔曼定理，$-\langle T\rangle/\mu=E_n/\mu$，所以
		\[
			\boxed{\langle T\rangle_n=-E_n
			=\frac{\mu e^4}{2\hbar^2n^2}}.
		\]
	\end{enumerate}
}

\newpage
\item 
（$30'$）粒子在二维无限深方势阱中运动，$V=\begin{cases}0,&0<x<a,0<y<a\\ \infty,&\text{其他}\end{cases}$。加上微扰$H'=\lambda xy$后，求基态和第一激发态能级的一级微扰修正。

\banswer{
	无扰本征态记为 $|n_x,n_y\rangle$，能量
	\[
		E_{n_xn_y}^{(0)}=\frac{\pi^2\hbar^2}{2ma^2}(n_x^2+n_y^2).
	\]
	基态 $|1,1\rangle$ 非简并，且任一一维箱中 $\langle n|x|n\rangle=a/2$，故
	\[
		\boxed{E_{11}=\frac{\pi^2\hbar^2}{ma^2}+\frac{\lambda a^2}{4}+O(\lambda^2)}.
	\]

	第一激发态由 $|1,2\rangle$、$|2,1\rangle$ 构成二重简并子空间。利用
	\[
		\langle1|x|2\rangle=\frac1a\int_0^a x
 \left(\cos\frac{\pi x}{a}-\cos\frac{3\pi x}{a}\right)dx
 =-\frac{16a}{9\pi^2},
	\]
	微扰矩阵为
	\[
		W=\lambda a^2\begin{pmatrix}
		\frac14&\frac{256}{81\pi^4}\\[2pt]
		\frac{256}{81\pi^4}&\frac14
		\end{pmatrix}.
	\]
	故归一化零级态 $(|1,2\rangle\pm|2,1\rangle)/\sqrt2$ 的一级能量修正为
	\[
		\boxed{E_\pm^{(1)}=\lambda a^2
		\left(\frac14\pm\frac{256}{81\pi^4}\right)},
	\]
	相应总能量为 $5\pi^2\hbar^2/(2ma^2)+E_\pm^{(1)}$。
}


\item 
（$30'$）设粒子所处的外场均匀但与时间有关。即$V=V(t)$，与坐标$\vec{r}$无关。试将体系的含时薛定鄂方程分离变量，求方程解$\psi(\vec{r},t)$的一般形式，并取$V(t)=V_0\cos(\omega t)$，以一维情况为例说明$V(t)$的影响是什么。

\banswer{
	取分离形式 $\psi(\vec r,t)=\phi(\vec r)T(t)$，代入
	\[
		i\hbar\frac{\partial\psi}{\partial t}
		=\left[-\frac{\hbar^2}{2m}\nabla^2+V(t)\right]\psi
	\]
	可令自由粒子空间部分满足
	\[
		-\frac{\hbar^2}{2m}\nabla^2\phi=E\phi,
	\]
	而时间部分满足
	\[
		T(t)=T(0)\exp\!\left[-\frac{i}{\hbar}
		\left(Et+\int_0^tV(t')dt'\right)\right].
	\]
	故一般解是自由粒子各本征态的叠加，每一项均多出相同因子
	$\exp[-(i/\hbar)\int_0^tV(t')dt']$。

	当 $V(t)=V_0\cos\omega t$ 时，一维平面波解为
	\[
		\psi_k(x,t)=C e^{ikx-i\hbar k^2t/(2m)}
		\exp\!\left[-i\frac{V_0}{\hbar\omega}\sin\omega t\right].
	\]
	因此空间均匀的 $V(t)$ 只给所有态附加共同的含时相位，不改变概率密度、动量分布或任何可观测跃迁概率。
}



\end{enumerate}

